\DocumentMetadata{
  pdfversion=2.0,
  pdfstandard=ua-2,
  lang=en-US,
  tagging=on,
  tagging-setup={math/setup=mathml-SE,tabsorder=structure}
}
\documentclass[11pt,letterpaper]{article}
\usepackage[margin=0.68in,includefoot]{geometry}
\usepackage{newcomputermodern}
\usepackage{amsmath,amscd,mathtools}
\usepackage{tikz,adjustbox}
\providecommand{\square}{\mdlgwhtsquare}
\usepackage[hidelinks]{hyperref}
\hypersetup{
  pdftitle={Analysis PhD Exam, September 2009},
  pdfauthor={University of Florida Department of Mathematics},
  pdfsubject={Graduate mathematics examination},
  pdfdisplaydoctitle=true,
  bookmarksopen=true
}
\setcounter{secnumdepth}{0}
\setlength{\parindent}{0pt}
\setlength{\parskip}{0.28em}
\setlength{\emergencystretch}{3em}
\setlength{\leftmargini}{2.15em}
\setlength{\leftmarginii}{2.65em}
\setlength{\leftmarginiii}{3em}
\pagestyle{empty}
\raggedbottom
\allowdisplaybreaks
\NewTaggingSocketPlug{block/list/label}{examproblem}{%
  \tagstructbegin{tag=Lbl}\tagstructend%
  \tagstructbegin{tag=\LBody}%
  \tagstructbegin{tag=\ProblemHeadingTag,title-o={\ProblemHeadingText}}%
  \tagmcbegin{tag=\ProblemHeadingTag,actualtext={\ProblemHeadingText}}%
  #2%
  \tagmcend%
  \tagstructend%
}
\newenvironment{examproblems}[2]{%
  \def\ProblemHeadingTag{#2}%
  \AssignTaggingSocketPlug{block/list/label}{examproblem}%
  \begin{enumerate}%
  \setcounter{enumi}{#1}%
  \renewcommand{\labelenumi}{\textbf{\theenumi.}}%
  \setlength{\topsep}{0.35em}%
  \setlength{\partopsep}{0pt}%
  \setlength{\itemsep}{0.48em}%
  \setlength{\parsep}{0.18em}%
}{\end{enumerate}}
\newenvironment{examsubparts}{%
  \AssignTaggingSocketPlug{block/list/label}{default}%
  \begin{enumerate}%
  \setlength{\topsep}{0.18em}%
  \setlength{\partopsep}{0pt}%
  \setlength{\itemsep}{0.16em}%
  \setlength{\parsep}{0.1em}%
}{\end{enumerate}}
\newenvironment{examinstructions}{%
  \par\begingroup\itshape
}{%
  \par\endgroup\vspace{0.18em}
}
\makeatletter
\renewcommand\subsection{\@startsection{subsection}{2}{0pt}%
  {-1.25ex plus -.25ex minus -.1ex}%
  {0.55ex plus .1ex}%
  {\normalfont\large\bfseries}}
\renewcommand{\@maketitle}{%
  \newpage\null\vskip -1em%
  {\footnotesize\bfseries
    \href{https://www.ufl.edu/}{University of Florida}, \href{https://math.ufl.edu/}{Department of Mathematics}\par}%
  \vskip -0.5em%
  \tagstructbegin{tag=H1,title={Analysis PhD Exam, September 2009}}%
  \tagpdfparaOff%
  \tagmcbegin{tag=H1}%
    {\LARGE\bfseries\href{https://gma.math.ufl.edu/exams/analysis-phd/}{Analysis PhD Exam}, September 2009\par}%
  \tagmcend%
  \tagpdfparaOn%
  \tagstructend%
  \vskip 0.25em%
  \tagmcbegin{artifact}%
    \hrule height 1pt%
  \tagmcend%
  \vskip 1em%
}
\makeatother
\title{Analysis PhD Exam, September 2009}
\author{}
\date{}
\begin{document}
\maketitle
\thispagestyle{empty}
\begin{examinstructions}
Do 6 out of the 7 problems.
\end{examinstructions}
\begin{examproblems}{0}{H2}
\def\ProblemHeadingText{Problem 1.}
\item
State and prove the Fubini Theorems.
\def\ProblemHeadingText{Problem 2.}
\item
State and prove the Vitali integral convergence theorem.
\def\ProblemHeadingText{Problem 3.}
\item
Let~\(\mu\) be a signed measure on the~\(\sigma\)-algebra \(\mathcal S\). Define positive and negative set for~\(\mu\). Prove that if \(E\in\mathcal S\), then \(\mu^+(E)=\sup\{\mu(F):F\subset E,\ F\in\mathcal S\}\).
\def\ProblemHeadingText{Problem 4.}
\item
Let~\(Y\) be an orthogonal set in a Hilbert space. Prove that \(\sum_{y\in Y}y\) converges if and only if \(\sum_{y\in Y}\lVert y\rVert^2<\infty\).
\def\ProblemHeadingText{Problem 5.}
\item
Let~\(X\) be a Banach space, \(A\subset X\). Suppose that \(\sup_{a\in A}|x^*a|<\infty\), for each \(x^*\in X^*\). Prove \(\sup_{a\in A}\lVert a\rVert<\infty\).
\def\ProblemHeadingText{Problem 6.}
\item
Let~\(\mu\) and~\(\nu\) be finite measures such that each is absolutely continuous with respect to the other. Prove \(\frac{d\mu}{d\nu}=\frac{1}{\left(\frac{d\nu}{d\mu}\right)}\) a.e. \(\mu\).
\def\ProblemHeadingText{Problem 7.}
\item
Let \(X,Y,Z\) be Banach spaces and let \(\mathcal F\) be a total family of continuous linear maps on~\(X\) to~\(Y\). Let \(T:Z\to X\) be a linear map such that \(fT\) is continuous for every \(f\in\mathcal F\). Prove~\(T\) is continuous. [Hint: Use the closed graph theorem.] Note: \(\mathcal F\) total means if \(f(x)=0\) for all \(f\in\mathcal F\), then \(x=0\).
\end{examproblems}
\end{document}
